\makeatletter

\pgfkeys{%
	/nx/.is family,%
	/nx/.style={
		/nx/.cd%
	},%
	%%% Cover Page Family %%%
	/nx/cover page/.is family,%
	nx cover page/.style={%
		nx, cover page/.cd%
	},%
	/nx/cover page/.cd,%
	% Default Values
	title/.default={Posix-Nexus},%
	author/.default={Canine-Table},%
	image/.default={img/posix-nexus.png},%
	link/.default={https://github.com/Canine-Table/posix-nexus},%
	abstract/.default={POSIX Nexus serves as a comprehensive cross-language reference hub that explores the implementation and behavior of POSIX-compliant functionality across a diverse set of programming environments. Built atop the foundational IEEE Portable Operating System Interface (POSIX) standards, this project emphasizes compatibility, portability, and interoperability between operating systems.},%
	Abstract/.default={Abstract},
	keywords/.default={%
		posix,%
		portability,%
		compatibility,%
		system calls,%
		apis,%
		threads,%
		shells,%
		utilities,%
		filesystems,%
		file handling%
		permissions,%
		signals,%
		nexus%
	},%
	subject/.default={%
		The Portable Operating System Interface%
	},%
	% Key Storage
	title/.store in=\nx@title,%
	author/.store in=\nx@author,%
	image/.store in=\nx@image,%
	abstract/.store in=\nx@abstract,%
	Abstract/.store in=\nx@Abstract,%
	link/.store in=\nx@link,%
	subject/.store in=\nx@subject,%
	keywords/.store in=\nx@keywords,%
	% Code
	ref/.code={%
		\pgfkeysifdefined{/nx/cover page/link}{}{%
			\pgfkeys{/nx/cover page/link}%
		}%
		\ifx#1\pgfkeysnovalue%
			\pgfkeysifdefined{/nx/cover page/ref}{}{%
				\def\nx@tmpa{\nx@link}%
			}%
		\else%
			\def\nx@tmpa{#1}%
		\fi%
		\pgfkeysgetvalue{/nx/cover page/ref}\nx@tmpb%
		\ifx\nx@tmpa\nx@tmpb\else%
			\pgfkeyssetvalue{/nx/cover page/ref}{\nx@tmpa}%
		\fi%
	},%
	%
	%
	/nx/refer/.is family,%
	nx refer/.style={%
		nx, refer/.cd%
	},%
	/nx/refer/.cd,%
	acr style/.default={long-short-desc},%
	gls style/.default={index},%
	bib style/.default={},%
	abr style/.default={long-short-sc},%
	acr style/.store in=\nx@acr,%
	gls style/.store in=\nx@gls,%
	bib style/.store in=\nx@bib,%
	abr style/.store in=\nx@abr,%
	%
	%
	%%% List Box Family %%%
	/nx/list box/.is family,%
	nx list box/.style={%
		nx, list box/.cd%
	},%
	/nx/list box/.cd,%
	% Code Defs
	location/.code 2 args={\nxMapMarkerDangerItem\nx@Topic{#1}{Location}{#2}},%
	module/.code 2 args={\nxBoxOpenObsoleteItem\nx@Topic{#1}{Module}{#2}},%
	parameters/.code 2 args={\nxCodeSuccessItem\nx@Topic{#1}{Parameters}{#2}},%
	class/.code 2 args={\nxBoxesObsoleteItem\nx@Topic{#1}{Class}{#2}},%
	type/.code 2 args={\nxToolboxDangerItem\nx@Topic{#1}{Type}{#2}},%
	summary/.code 2 args={\nxRocketInfoItem\nx@Topic{#1}{Summary}{#2}},%
	behavior/.code 2 args={\nxSearchInfoItem\nx@Topic{#1}{Behavior}{#2}},%
	dependencies/.code 2 args={\nxLinkInfoItem\nx@Topic{#1}{Dependencies}{#2}},%
	reference/.code 2 args={\nxThumbTackDangerItem\nx@Topic{#1}{Reference}{#2}},%
	availability/.code 2 args={\nxBrainExperimentalItem\nx@Topic{#1}{Availability}{#2}},%
	true/.code 2 args={\nxCheckSuccessItem\nx@Topic{#1}{True}{#2}},%
	false/.code 2 args={\nxTimesDangerItem\nx@Topic{#1}{False}{#2}},%
	/nx/shapes/.is family,%
	nx shapes/.style={%
		nx, shapes/.cd%
	},%
	/nx/shapes/.cd,%
	byX/.initial=-9pt,%
	byY/.initial=-9pt,%
	%
	%
	/nx/triangle/.is family,%
	nx triangle/.style={%
		nx, triangle/.cd%
	},%
	/nx/triangle/.cd,%
	al/.default=A,%
	bl/.default=B,%
	cl/.default=C,%
	as/.default={-2.5,2.5},%
	bs/.default={2.5,2.5},%
	cs/.default={-2.5,-2.5},%
	al/.store in=\nx@AL,%
	bl/.store in=\nx@BL,%
	cl/.store in=\nx@CL,%
	as/.store in=\nx@AS,%
	bs/.store in=\nx@BS,%
	cs/.store in=\nx@CS,%
	%
	%
	/nx/crowsfoot/.is family,%
	nx crowsfoot/.style={%
		nx, crowsfoot/.cd%
	},%
	/nx/crowsfoot/.cd,%
	shaftoffset/.initial=0.25,%
	tipoffset/.initial=0.35,%
	barA/.initial=0.6,%
	barB/.initial=0.3,%
	circleradius/.initial=0.25,%
	circleoffset/.initial=0.03,%
	circlestart/.initial=0pt,%
	footFork/.initial=0.35,%
	footSpread/.initial=0.3,%
	backend/.initial=0.75,%
	visualbackend/.initial=0.5,%
	lineend/.initial=0.25,%
	tipend/.initial=0.25,%
	tipoffset/.initial=0pt%
}

\pgfdeclarearrow{
	name = nxMM,
	parameters = {
		\the\pgfarrowlength
	},
	setup code = {
		% The different end values:
		\pgfarrowssettipend{.25\pgfarrowlength}
		\pgfarrowssetlineend{-.25\pgfarrowlength}
		\pgfarrowssetvisualbackend{-.5\pgfarrowlength}
		\pgfarrowssetbackend{-.75\pgfarrowlength}
		% The hull
		\pgfarrowshullpoint{.25\pgfarrowlength}{0pt}
		\pgfarrowshullpoint{-.75\pgfarrowlength}{.5\pgfarrowlength}
		\pgfarrowshullpoint{-.75\pgfarrowlength}{-.5\pgfarrowlength}
		% Saves: Only the length:
		\pgfarrowssavethe\pgfarrowlength
	},
	drawing code = {
		\pgfpathmoveto{\pgfqpoint{.35\pgfarrowlength}{0pt}}
		\pgfpathlineto{\pgfqpoint{\pgfarrowlength}{.3\pgfarrowlength}}
		\pgfpathmoveto{\pgfqpoint{.35\pgfarrowlength}{0pt}}
		\pgfpathlineto{\pgfqpoint{\pgfarrowlength}{-.3\pgfarrowlength}}
		\pgfpathmoveto{\pgfqpoint{-.25\pgfarrowlength}{0pt}}
		\pgfpathlineto{\pgfqpoint{\pgfarrowlength}{0pt}}
		\pgfpathmoveto{\pgfqpoint{.25\pgfarrowlength}{.3\pgfarrowlength}}
		\pgfpathlineto{\pgfqpoint{.25\pgfarrowlength}{-.3\pgfarrowlength}}
		\pgfpathclose
		\pgfusepathqstroke
	},
	defaults = {
		length = 5mm
	}
}

\pgfdeclarearrow{
	name = nxMO,
	parameters = {
		\the\pgfarrowlength
	},
	setup code = {
		% The different end values:
		\pgfarrowssettipend{.25\pgfarrowlength}
		\pgfarrowssetlineend{-.25\pgfarrowlength}
		\pgfarrowssetvisualbackend{-.5\pgfarrowlength}
		\pgfarrowssetbackend{-.75\pgfarrowlength}
		% The hull
		\pgfarrowshullpoint{.25\pgfarrowlength}{0pt}
		\pgfarrowshullpoint{-.75\pgfarrowlength}{.5\pgfarrowlength}
		\pgfarrowshullpoint{-.75\pgfarrowlength}{-.5\pgfarrowlength}
		% Saves: Only the length:
		\pgfarrowssavethe\pgfarrowlength
	},
	drawing code = {
		\pgfpathmoveto{\pgfqpoint{-.25\pgfarrowlength}{0pt}}
		\pgfpathlineto{\pgfqpoint{\pgfarrowlength}{0pt}}
		\pgfpathmoveto{\pgfqpoint{.6\pgfarrowlength}{-.3\pgfarrowlength}}
		\pgfpathlineto{\pgfqpoint{.6\pgfarrowlength}{.3\pgfarrowlength}}
		\pgfpathmoveto{\pgfqpoint{.3\pgfarrowlength}{-.3\pgfarrowlength}}
		\pgfpathlineto{\pgfqpoint{.3\pgfarrowlength}{.3\pgfarrowlength}}
		\pgfpathclose
		\pgfusepathqstroke
	},
	defaults = {
		length = 5mm
	}
}

\pgfdeclarearrow{
	name = nxOO,
	parameters = {
		\the\pgfarrowlength
	},
	setup code = {
		% The different end values:
		\pgfarrowssettipend{.25\pgfarrowlength}
		\pgfarrowssetlineend{-.25\pgfarrowlength}
		\pgfarrowssetvisualbackend{-.5\pgfarrowlength}
		\pgfarrowssetbackend{-.75\pgfarrowlength}
		% The hull
		\pgfarrowshullpoint{.25\pgfarrowlength}{0pt}
		\pgfarrowshullpoint{-.75\pgfarrowlength}{.5\pgfarrowlength}
		\pgfarrowshullpoint{-.75\pgfarrowlength}{-.5\pgfarrowlength}
		% Saves: Only the length:
		\pgfarrowssavethe\pgfarrowlength
	},
	drawing code = {
		\pgfpathmoveto{\pgfqpoint{.6\pgfarrowlength}{.3\pgfarrowlength}}
		\pgfpathlineto{\pgfqpoint{.6\pgfarrowlength}{-.3\pgfarrowlength}}
		\pgfpathmoveto{\pgfqpoint{-.25\pgfarrowlength}{0pt}}
		\pgfpathlineto{\pgfqpoint{\pgfarrowlength}{0pt}}
		\pgfpathclose
		\pgfusepathqstroke
		\pgfpathcircle{\pgfpoint{.03\pgfarrowlength}{0pt}}{.25\pgfarrowlength}
		\pgfpathclose
		\pgfusepathqfill
	},
	defaults = {
		length = 5mm
	}
}

\pgfdeclarearrow{
	name = nxOM,
	parameters = {
		\the\pgfarrowlength
	},
	setup code = {
		% The different end values:
		\pgfarrowssettipend{.25\pgfarrowlength}
		\pgfarrowssetlineend{-.25\pgfarrowlength}
		\pgfarrowssetvisualbackend{-.5\pgfarrowlength}
		\pgfarrowssetbackend{-.75\pgfarrowlength}
		% The hull
		\pgfarrowshullpoint{.25\pgfarrowlength}{0pt}
		\pgfarrowshullpoint{-.75\pgfarrowlength}{.5\pgfarrowlength}
		\pgfarrowshullpoint{-.75\pgfarrowlength}{-.5\pgfarrowlength}
		% Saves: Only the length:
		\pgfarrowssavethe\pgfarrowlength
	},
	drawing code = {
		\pgfpathmoveto{\pgfqpoint{.35\pgfarrowlength}{0pt}}
		\pgfpathlineto{\pgfqpoint{\pgfarrowlength}{.3\pgfarrowlength}}
		\pgfpathmoveto{\pgfqpoint{.35\pgfarrowlength}{0pt}}
		\pgfpathlineto{\pgfqpoint{\pgfarrowlength}{-.3\pgfarrowlength}}
		\pgfpathmoveto{\pgfqpoint{-.25\pgfarrowlength}{0pt}}
		\pgfpathlineto{\pgfqpoint{\pgfarrowlength}{0pt}}
		\pgfpathclose
		\pgfusepathqstroke
		\pgfpathcircle{\pgfpoint{.03\pgfarrowlength}{0pt}}{.25\pgfarrowlength}
		\pgfpathclose
		\pgfusepathqfill
	},
	defaults = {
		length = 5mm
	}
}

\def\nxPlaceLabel#1#2{%
	\path (#1);%
	\pgfgetlastxy{\xcoord}{\ycoord}%
	\pgfmathsetmacro{\xsign}{\xcoord < 0 ? "left" : "right"}%
	\pgfmathsetmacro{\ysign}{\ycoord < 0 ? "below" : "above"}%
	\node[anchor=\ysign \xsign] at (#1) {#2};%
}

\pgfdeclareshape{simple shape}{%
	\savedanchor{\center}{%
		\pgfpointorigin}
	\anchor{center}{\center}
	\savedanchor{\anchorfoo}{%
		\pgf@x=1cm
		\pgf@y=0cm}
	\deferredanchor{anchor \foo}{\anchorfoo}}

\def\polartransformation{%
	% \pgf@x will contain the radius
	% \pgf@y will contain the distance
	\pgfmathsincos@{\pgf@sys@tonumber\pgf@x}%
	% pgfmathresultx is now the cosine of radius and
	% pgfmathresulty is the sine of radius
	\pgf@x=\pgfmathresultx\pgf@y%
	\pgf@y=\pgfmathresulty\pgf@y%
}

\pgfdeclareverticalshading{shading1}{100bp}{%
	color(0bp)=(red);%
	color(25bp)=(green);%
	color(75bp)=(blue);%
	color(100bp)=(black)%
}

\pgfdeclareverticalshading{shading2}{20pt}{%
	color(0pt)=(red);%
	color(20pt)=(blue)%
}

\pgfdeclareradialshading{shading3}{\pgfpoint{-10bp}{10bp}}{%
	color(0bp)=(red!15!white);%
	color(9bp)=(red!75!white);%
	color(18bp)=(red!70!black);%
	color(25bp)=(red!50!black);%
	color(50bp)=(black)%
}%

\pgfdeclareradialshading{shading4}{\pgfpointorigin}{
	color(0mm)=(pgftransparent!0);
	color(5mm)=(pgftransparent!0);
	color(8mm)=(pgftransparent!100);
	color(15mm)=(pgftransparent!100)
}

\pgfdeclarehorizontalshading{shading5}{100bp}{
	color(0pt)=(transparent!0);
	color(25bp)=(transparent!0);
	color(75bp)=(transparent!100);
	color(100bp)=(transparent!100)
}

\pgfdeclareverticalshading{shading6}{4cm}{
	rgb(0cm)=(1,0,0);
	rgb(1.5cm)=(0,1,0);
	rgb(2cm)=(0,0,1)
}


\pgfdeclareradialshading{shading7}{\pgfpoint{0.5cm}{0.5cm}}{%
	rgb(0cm)=(0.9,0,0);%
	rgb(0.7cm)=(0.7,0,0);%
	rgb(1cm)=(0.5,0,0);%
	rgb(1.05cm)=(1,1,1)%
}%

\pgfdeclarefunctionalshading[black]{portabletwospots}{\pgfpointorigin}{\pgfpoint{3.5cm}{3.5cm}}{}{
	2 copy
	45 sub dup mul exch
	40 sub dup mul 0.5 mul add sqrt
	dup mul neg 1.0005 exch exp 1.0 exch sub
	3 1 roll
	70 sub dup mul .5 mul exch
	70 sub dup mul add sqrt
	dup mul neg 1.002 exch exp 1.0 exch sub
	2.0 3 1 roll
	\ifpgfshadingmodelcmyk
		\pgffuncshadingrgbtocmyk
	\fi
	\ifpgfshadingmodelgray
		\pgffuncshadingrgbtogray
	\fi
}

\pgfdeclarelayer{background layer}
\pgfdeclarelayer{foreground layer}

\pgfsetlayers{background layer,main,foreground layer}

\makeatother

\pgfooclass{NxLuaObject}{
	\attribute nxClass;
	\attribute nxInstance;
	\method NxLuaObject(#1) {
		\pgfooset{nxClass}{#1}%
		\pgfoothis.get id(\nxInstance)%
		\pgfooset{nxInstance}{nex.tex.oop.oo["\pgfoovalueof{nxClass}"].instance["\nxInstance"]}%
		\directlua{nex.tex.oop:new("\pgfoovalueof{nxClass}", "\nxInstance")}%
	}

	\method setProperties(#1) {
		\directlua{\pgfoovalueof{nxInstance}:setProperties({#1})}
	}
}

\pgfooclass(NxLuaObject){NxTikzPGF}{
	\method NxTikzPGF() {
		\pgfoothis.NxLuaObject.NxLuaObject(NxTikzPGF)
	}

	\method triangle() {
		\directlua{\pgfoovalueof{nxInstance}:pgfTriangle()}
	}

	\method rectangle() {
		\directlua{\pgfoovalueof{nxInstance}:pgfRectangle()}
	}

	\method circle() {
		\directlua{\pgfoovalueof{nxInstance}:pgfCircle()}
	}
}

